\subsection{}
\label{XIII.2.10}
If $U$ is a connected scheme, $a$ a geometric point
of~$U$, $\LL$~a set of prime numbers, we denote by
\begin{equation*}
\label{eq:XIII.2.10.0}
\tag{2.10.0} \pi_1^{\LL}(U,a)
\end{equation*}
the projective limit of the finite quotients of~$\pi_1(U,a)$ whose
orders have all their prime factors in~$\LL$.

We shall define specialization morphisms for the
fundamental group, generalizing X.\Ref{X.2}.

Let $g\colon U\to S$ be a coherent morphism with geometrically
connected fibers (\resp a morphism of the form $g=
fi$, where $f\colon X \to S$ is a proper morphism of
finite presentation and where $i\colon U\to X$ is an
open immersion such that $U$ is the complement in $X$
of a divisor with normal crossings relative to~$S$
(\cf \Ref{XIII.2.8})). Let ${\LL}$ be a set of prime numbers and
suppose, in the non-resp.\ case, that, for every finite constant
${\LL}$-group $C$, $(C_U,g)$ is cohomologically proper
relative to~$S$ in dimension $\leq 1$. Let $\overline{s}_1
\to \overline{s}_2$ be a specialization morphism of
geometric points of~$S$, $\overline{S}$ the strict
localization of~$S$ at $\overline{s}_2$,
$\overline{U}=U\times_S\overline{S}$. We have a commutative diagram
$$
\xymatrix{ U_{\bar{s}_1} \ar[r]^{h_1} \ar[d] & \bar{U}
\ar[d]^{\bar{g}} & U_{\bar{s}_1} \ar[l]_{h_2} \ar[d] \\ \bar{s}_1
\ar[r] & \bar{S} & \,\bar{s}_2\,. \ar[l] }
$$
If $a_1$ is a geometric point of~$U_{\bar{s}_1},\ a_2$ a
geometric point of~$U_{\bar{s}_2}$, the morphisms $h_1$ and
$h_2$ define canonical morphisms
% original p. 391
\begin{align*}
\pi_1\colon \pi_1^{\LL}(U_{\bar{s}_1},a_1)&\to
\pi_1^{\LL}(\bar{U},a_1) & \pi_2\colon
\pi_1^{\LL}(U_{\bar{s}_2},a_2)&\to \pi_1^{\LL}(\bar{U},a_2)\\
(\text{\resp~} \pi_1\colon \pi_1^\tame (U_{\bar{s}_1},a_1)&\to
\pi_1^\tame (\bar{U},a_1) & \pi_2\colon
\pi_1^\tame(U_{\bar{s}_2},a_2)&\to \pi_1^\tame(\bar{U},a_2))
\end{align*}
(V~\Ref{V.7} and~\Ref{eq:XIII.2.1.5.2}). The hypotheses of
cohomological properness (\resp \Ref{XIII.2.8}) prove that $\pi_2$
is an isomorphism. If one chooses a class of paths from~$a_1$
to $a_2$, one obtains an isomorphism
$$
\pi_{12}\colon \pi_1^{\LL}(\bar{U},a_1)\isomto
\pi_1^{\LL}(\bar{U},a_2) \quad (\text{\resp~} \pi_{12}\colon
\pi_1^\tame (\bar{U},a_1)\isomto \pi_1^\tame(\bar{U},a_2))
$$
whence a morphism $\pi=\pi_2^{-1}\pi_{12}\pi_1$
$$
\pi\colon \pi_1^{\LL}(U_{\bar{s}_1},a_1)\to
\pi_1^{\LL}(U_{\bar{s}_2},a_2) \quad (\text{\resp~}\pi\colon
\pi_1^\tame (U_{\bar{s}_1},a_1)\to \pi_1^\tame
(U_{\bar{s}_2},a_2)).
$$
Changing the class of paths from~$a_1$ to $a_2$ amounts to
modifying $\pi$ by an inner automorphism of
$\pi_1^{\LL}(X_{\bar{s}_2},a_2)$ (\resp of
$\pi_1^\tame(X_{\bar{s}_2},a_2)$). We call \emph{specialization
morphism for the fundamental group}
associated with the morphism $\bar{s}_1\to \bar{s}_2$, and we denote
simply by
$$
\pi\colon \pi_1^{\LL}(X_{\bar{s}_1})\to
\pi_1^{\LL}(X_{\bar{s}_2}) \quad \text{(\resp~} \pi\colon
\pi_1^{\tame}(X_{\bar{s}_1})\to \pi_1^{\tame}(X_{\bar{s}_2}))
$$
any one of the morphisms defined above.

\begin{lemma}
\label{XIII.2.11}
Let $f\colon X\to S$ be a proper morphism of
finite presentation, $D$ a divisor on~$X$ with normal
crossings relative to~$S$, $Y=\Supp D$, $U=X-Y$, $i\colon
U\to X$ the canonical morphism, $\bar{s}_1\to
\bar{s}_2$ a specialization morphism of geometric points
of~$S$, $y_1$ a geometric point of
$Y_{\bar{s}_1}$, $y_2$ a geometric point of
$Y_{\bar{s}_2}$, such that the projection $z_1$ of~$y_1$ on~$X$ is a
generization of the projection $z_2$ of~$y_2$. Let $I_{y_1}$
be an inertia subgroup of~$\pi_1^{\tame}(\bar{U}_{\bar{s}_1})$ at
$y_1.$ Then the image of~$I_{y_{1}}$ under the
specialization morphism
$$
\pi\colon \pi_1^{\tame}(U_{\bar{s}_1})\to
\pi_1^{\tame}(U_{\bar{s}_2})
$$
is
% original p. 392
an inertia subgroup of~$\pi_1^{\tame}(U_{\overline{s}_2})$
at~$y_2$.
\end{lemma}

Indeed, let $\overline{X}$ (\resp $\widetilde{X}$) be the strict
localization of~$X$ at $y_2$ (\resp at $y_1$), $\overline{U}=U \times_X
\overline{X}$ (\resp $\widetilde{U}=U\times_X \widetilde{X}$). We have a
canonical morphism $\widetilde{U}\to \overline{U}$, and it
follows from~\Ref{XIII.1.10} that we have a commutative diagram
$$
\xymatrix{ \pi_1^{\tame} \left(\widetilde{U}_{\overline{s}_1} \right)
\ar[r]^{\pi '} \ar[d] & \pi_1^{\tame} \left(
\overline{U}_{\overline{s}_2} \right) \ar[d] \\ \pi_1^{\tame} \left(
U_{\overline{s}_1} \right) \ar[r]^{\pi } & \pi_1^{\tame} \left(
U_{\overline{s}_2} \right) }
$$
where $\pi'$ is the composite of the canonical morphism
$\pi_1^{\tame}(\widetilde{U}_ {\overline{s}_1})\to
\pi_1^{\tame}(\overline{U}_{\overline{s}_1})$ and of the
specialization morphism. Since
$\pi_1^{\tame}(\widetilde{U}_{\overline{s}_1})$
(\resp $\pi_1^{\tame}(\overline{U}_{\overline{s}_1})$)
is an inertia group of
$\pi_1^{\tame}(U_{\overline{s}_1})$ at~$y_1$ and (\resp of~$\pi_1^{\tame}(U_{\overline{s}_2})$ at~$y_2$), it suffices to
prove that $\pi'$ is surjective. But this follows from
the expression obtained in~\Ref{XIII.5.6}.

\begin{corollary}
\label{XIII.2.12}
Let $X$ be a proper and smooth connected curve of genus $g$ over a
separably closed field $k$ of characteristic $p\geq
0$. Let $U$ be the open set obtained by removing from $X$ $n$ distinct
closed points $a_1,\dots,a_n$. Then the tamely ramified
fundamental group $\pi_1^{\tame}(U)$ \eqref{XIII.2.1.3}
can be generated by $2g+n$ elements $x_i,
y_i,\sigma_j$, with $1\leq i\leq g$, $1\leq j\leq n$, such that
$\sigma_j$ is a generator of an inertia group
corresponding to $a_j$, and such that we have the relation
\begin{equation*}
\label{eq:XIII.2.12.*}
\tag{$*$} \prod_{1\leq i\leq
g}(x_iy_ix^{-1}_iy^{-1}_i)\cdot\prod_{1\leq j\leq n}\sigma_j=1.
\end{equation*}
For every finite group $G$ \emph{of order prime to} $p$,
generated by elements $\overline{x}_i,
\overline{y}_i,\overline{\sigma}_j$ satisfying the relation
\eqref{eq:XIII.2.12.*}, there exists an \'etale covering of~$U$,
with group $G$, corresponding to a homomorphism
$\pi_1^{\tame}(U)\to G$ which
% original p. 393
sends $x_i, y_i,\sigma_j$ to~$\overline{x}_i,
\overline{y}_i,\overline{\sigma}_j$ respectively. In other
words, if $p'$ denotes the set of prime numbers distinct
from~$p$, $\pi_1^{p'}(U)$ is the pro-$p'$-group generated by the
generators $x_i, y_i,\sigma_j$ bound by the single relation~\eqref{eq:XIII.2.12.*}.
\end{corollary}

\subsubsection*{Proof}
We may assume $k$ algebraically closed. Suppose first $k$ of
characteristic zero. There then exists an algebraically closed
subextension $k'$ of~$k$, of finite transcendence degree
over~$\QQ$, such that $X$ comes by extension of scalars
from a proper and smooth curve $X'$ defined over~$k'$, and we may
assume that the points $a_1,\dots,a_n$ come from rational points
$a'_1,\dots,a'_n$ of~$X'$. Since $k'$ is of finite transcendence
degree over~$\QQ$, we can find an embedding of~$k'$ into the field
of complex numbers $\CC$; let $\widetilde{U}=U'\times_{k'}\CC.$
Let $k''$ be an algebraically closed extension of~$k'$ such that
we have $k'$-morphisms of~$k$ and of~$\CC$ into $k''$. If
$g'\colon U'\to k'$ is the structural morphism, and if
$\cal{F}$ is a finite constant sheaf of groups on~$U''$, it
follows from~\Ref{XIII.2.9} that the specialization morphisms
$$
(\R^1g'_*\cal{F}')_{\CC}\to(\R^1g'_*\cal{F}')_{k''}
\from(\R^1g'_*\cal{F}')_k
$$
are isomorphisms. In terms of fundamental groups, this shows
that we have an isomorphism, defined up to inner
automorphism
$$
\pi_1(U)\to\pi_1(\widetilde{U}),
$$
and it is clear that this isomorphism transforms an inertia group
relative to a point of~$X'-U'$ into an inertia group relative to the
same point. We may therefore assume that $k=\CC$. In this
last case it follows from Riemann's existence theorem
(XII~\Ref{XII.5.2}) that the fundamental group $\pi_1(U)$ is none other
than the completion, for the topology of subgroups of finite
index, of the fundamental group of the analytic space associated with~$U$. Now the latter can be computed by transcendental means \cite[ch\ptbl 7
\S 47]{XIII.3};
% original p. 394
it can be generated by $2g+n$ elements $x_i,
y_i,\sigma_j$ such that $\sigma_j$ is the image of a generator
of the local fundamental group $\pi_1(D_j)$ of a small disk centered
at $a_j$, \ie a generator of an inertia group
corresponding to the point $a_j$, these elements satisfying
the single relation \eqref{eq:XIII.2.12.*}.

If now $k$ is of
characteristic $p>0$, we can find a complete discrete valuation
ring $A$, with residue field $k$, with field of
fractions $K$ of characteristic zero, and a connected
scheme $X_1$, proper and smooth over~$S=\Spec A$, such that we have
$X_1\times_S \Spec k\simeq X$ (III~\Ref{III.7.4}). The points $a_j$ then
lift to sections $s_j$ of~$X_1$ over~$S$;
let $Y_{1j}$ be the reduced closed subscheme with
underlying space $s_j(S)$, $Y_1$ the union of the $Y_{1j}$,
$U_1=X_1-Y_1$, $g_1\colon U_1\to S$ the structural
morphism. Let $\overline{K}$ be an algebraically closed extension
of~$K$, $\overline{U}= U_1\times_S\overline{K}$. If $C$ is a
finite constant group, it follows from~\Ref{XIII.2.8} that the
specialization morphism
$$
(\R^1g_{1*}C_{U_1})_k\to(\R^1g_{1*}C_{U_1})_{\overline{K}}
$$
is injective and even bijective if $C$ is of order prime to
$p$. Now this means, in terms of fundamental groups, that the
specialization morphism \eqref{XIII.1.10}
$$
\pi \colon \pi_1(\overline{U})\to\pi^{\tame}_1(U)
$$
is surjective, and that the specialization morphism
$$
\pi^{p'}_1(\overline{U})\to\pi^{p'}_1(U)
$$
is bijective. Finally, if $x_i, y_i,\sigma_j$ are
generators of~$\pi_1(\overline{U})$ such that $\sigma_j$ is
a generator of an inertia group corresponding to the point
$b_j=Y_{1j}(\overline{K})$ of~$\overline{X}$, then,
by~\Ref{XIII.1.11}, $\pi(\sigma_j)$ is a generator
of an inertia group corresponding to $a_j$, which completes the
proof.

\begin{remarkMR}
\label{XIII.2.13}
Let $k$ be an algebraically closed field of characteristic $p>0$, $X$~a proper smooth algebraic curve over~$k$, connected, of genus $g$, $U$~an affine open set of~$X$, the complement of~$r\geq1$ rational points
of~$X$. We have at our disposal the fundamental group $\pi_1(U)$, its quotient
$\pi_1^{\tame}(U)$ which classifies the finite \'etale coverings of~$U$
tamely ramified at the points of~$X-U$, and the quotient
$\pi_1^{p'}(U)$ of~$\pi_1^{\tame}(U)$ which classifies the Galois
\'etale coverings of~$U$ with Galois group of order prime to~$p$.

In (Coverings of algebraic curves, Amer.\ J.~Math.\ \textbf{79}
(1957), p\ptbl825--856) S\ptbl Abhyankar formulated a certain number of conjectures on the structure of the finite groups that are Galois groups of finite
\'etale coverings of~$U$.

The conjectures concerning the finite groups that are Galois groups
of connected \'etale coverings of~$U$ of order prime to~$p$
are proved and made more precise in Corollary \Ref{XIII.2.12}. Before tackling the finite quotients of~$\pi_1(U)$, let us begin by giving some
indications about the ``size'' of this group. Let $X = \Spec(A)$. It is
known that $\Hom_{\mathrm{cont}}(\pi_1(U), \ZZ/p\ZZ)$ is described by
Artin-Schreier theory (cor\ptbl XI~\Ref{XI.6.9}). This group is
isomorphic to $A/\wp(A)$ where $\wp$ is the map from~$A$ to $A$ which
sends $a$ to~$a^p-a$. Suppose first that $U$ is the affine
line $\AA^1$, with ring $k[T]$. Let~$E$ be the set of
elements of~$k[T]$ of the form $\sum_i a_iT^i$, with $a_i = 0$
when $p$ divides~$i$. It is immediate that the composite map
$E\to k[T]\to k[T]/\wp k[T]$ is bijective. Consequently the coefficients
$a_i$, $(i,p)=1$, behave like coordinates of a space which
parametrizes the cyclic coverings of degree
$p$ of the affine line. In particular, one deduces that $\pi_1(\AA^1)$ is not
topologically of finite type and that, if one makes an extension
$k\to k'$ of algebraically closed fields, the natural morphism $\pi_1(\AA^1\times_kk')\to\pi_1(\AA^1)$, which is surjective, is not bijective,
contrary to the proper case. In the general case, the curve $U$ is
realized as a scheme finite over the affine line and one concludes that the
same phenomena occur for $\pi_1(U)$ (and moreover, more
generally, for $\pi_1(V)$ for every connected affine $k$-scheme $V$
of finite type, of dimension $\geq1$).

This being so, if $G$ is a finite group, let us denote by $G^{(p')}$ the largest
quotient group of~$G$ of order prime to~$p$. For a finite group
$G$ to be a topological quotient of~$\pi_1(U)$, it is necessary that $G^{(p')}$
be a topological quotient of~$\pi_1^{p'}(U)$, a condition which
one knows in principle how to answer thanks to Corollary \Ref{XIII.2.12}. In the
aforementioned article, Abhyankar conjectures that this necessary condition is
also sufficient. This conjecture was proved by M\ptbl Raynaud
in the case of the affine line and by D\ptbl Harbater in the general case
(Invent. Math.\ \textbf{116} (1994), p\ptbl425--462 and \textbf{117}, p\ptbl1--25). For example,
in the case of the affine line $\pi_1^{p'}(\AA^1)=1$ and one
concludes that a finite group $G$ is the Galois group of a connected
\'etale covering of~$\AA^1$ if and only if $G^{(p')} = 1$,
that is, if and only if $G$ is generated by its
Sylow $p$-subgroups. Thus every finite simple group of order a
multiple of~$p$ will do.
\end{remarkMR}


\section{Cohomological properness and generic local acyclicity}
\label{XIII.3}
% original p. 395

\begin{theorem}
\label{XIII.3.1}
Let $S$ be an irreducible scheme with generic
point~$s$, $X$ and $Y$ two $S$-schemes of
finite presentation, $f\colon X \to Y$ an $S$-morphism. For every
$S$-scheme~$S'$, we denote by $Y'$, $X'$, etc.\ the inverse image of~$Y$,
$X$, etc.\ under the morphism $S' \to S$. We have the following
properties:

\begin{enumerate}
\item[1)]
{\upshape a)} One can find a nonempty open set~$S'$ of~$S$ such that, for
every finite constant sheaf of sets $F'$ on~$X'$, $f'_*F'$ is
constructible, and $(F',f')$ is cohomologically proper
relative to~$S'$ in dimension $\le 0$.

{\upshape b)} Let $F$ be a constructible sheaf of sets on~$X$. Then
one can find a nonempty open set~$S'$ of~$S$ (depending on~$F$)
such that $f'_*F'$ is constructible and $(F',f')$ is
cohomologically proper relative to~$S'$ in dimension $\le 0$.

\item[2)] Suppose that the schemes of finite type of dimension $\le
\dim X_s$ over an algebraic closure $\overline{k}$ of~$\kres(s)$
are strongly desingularizable \textup{(SGA~5 I~3.1.5)}. Then we have
moreover the following properties:

{\upshape a)} One can find a nonempty open set~$S'$ of~$S$ such that, for
every finite constant sheaf of groups $F'$ on~$X'$, of order prime
to the residue characteristics of~$S$, if $\Phi'$ is the
stack of torsors under $F'$, $f'_*\Phi'$ is $1$-constructible, and
$(F',f')$ is cohomologically proper relative to~$S'$ in
dimension $\le 1$.

{\upshape b)} Let $\LL$ be the set of prime numbers distinct from the
residue characteristics of~$S$ and $\Phi$ a
$1$-constructible ind-$\LL$-stack on~$X$ \eqref{XIII.0}, such that,
for every scheme~$X_1$ \'etale over~$X$ and for every pair
of objects $x, x_1$ of~$\Phi_{X_1}$, the sheaf $\SheafHom
_{X_1}(x,x_1)$ is constructible. Suppose moreover $S$ locally
noetherian. Then one can find a nonempty open set~$S'$ of~$S$
such that $f'_*\Phi'$ is $1$-constructible, that, for every pair
of objects $y, y_1$ of a fiber $(f'_*\Phi')_{Y_1}$,
$\SheafHom_{Y_1}(y,y_1)$ is constructible,
% original p. 396
and that $(\Phi',f')$ is cohomologically proper relative to~$S'$ in dimension $\le 1$.
\end{enumerate}
\end{theorem}

\subsubsection*{Proof}
We may assume $S$ affine; by SGA~4 VIII~1.1, we may
assume $S$ integral; finally, by passage to the limit, we may
assume that $S$ is the spectrum of an algebra of finite type over~$\ZZ$; in particular $S$ is then noetherian. Since the question
is local on~$Y$, we may assume $Y$ affine. Moreover, to prove the
theorem, it suffices to do so after a
finite extension $S'\to S$, where $S'$ is an integral scheme
and where $S'\to S$ is a composite of \'etale morphisms and of
finite radicial surjective morphisms.

\Subsection*{1) Case of constant sheaves of sets}

1) 1. Reduction to the case where $X$ is normal over~$S$. Let
$X_{1\overline{s}}$ be the normalization of~$(X_{\overline{s}})_{\red}$;
after restricting $S$ to a nonempty open set and making
a radicial extension of~$S$ if necessary, we may assume that
$X_{1\overline{s}}$ comes from a scheme $X_1$ normal over~$S$, and
that the morphism $X_{1\overline{s}}\to X_{\overline{s}}$ comes from a
finite surjective morphism $p\colon X_1\to X$ (EGA~IV 8.8.2 and
9.6.1). Suppose the theorem proved for
$fp$. After restricting $S$ to an open set if necessary, we may assume
that, for every constant sheaf of sets $F$ on~$X$, $(p^*F,fp)$
is cohomologically proper relative to~$S$ in dimension $\le
0$ and that $f_*p_*(p^*F)$ is
constructible. By~\Ref{XIII.1.9}, $(p_*p^*F,f)$ is then
cohomologically proper relative to~$S$ in dimension $\le
0$. The morphism
$$
F \to p_* p^* F=G
$$
is a monomorphism. It already follows from this, $f_*F$ being
a subsheaf of~$f_*G$, that $f_*F$ is constructible (SGA~4
IX~2.9~(ii)) and that $(F,f)$ is cohomologically proper relative
to~$S$ in dimension $\le -1$.

Let $X_2=X_1 \times_{X} X_1$, $q \colon X_2 \to X$ the
canonical morphism. By \Ref{XIII.1.11}~1),
% original p. 397
we have an exact sequence
$$
\xymatrix@C=.5cm{F \ar[r] & G \ar@<2pt>[r]\ar@<-2pt>[r] & q_*q^*F }.
$$

By what has just been proved, applied to
$fq$ instead of~$f$, we may assume, after restricting $S$
to a nonempty open set if necessary, that, for every constant sheaf of sets
$F$ on~$X$, $(q^*F,fq)$ is cohomologically proper relative
to~$S$ in dimension $\le -1$, hence that $(q_*q^*F,f)$ is
cohomologically proper relative to~$S$ in dimension $\le
-1$. It then follows from \Ref{XIII.1.13}~1) that $(F,f)$ is
cohomologically proper relative to~$S$ in dimension $\le 0$.

1) 2. Reduction to the case where $X$ is normal affine over~$S$.

Let $U_s$ be an affine open set of~$X_s$ dense in $X_s$. After
restricting $S$ to a nonempty open set if necessary, we may assume that $U_s
\to X_s$ lifts to an open immersion $i\colon U\to X$,
schematically dominant relative to~$S$ (EGA~IV
8.9.1). Since the morphism $X \to S$ is normal, we have by
SGA~2 XIV~1.18:
$$
\prof \et_{S- U} (X) \ge 2 \quoi ;
$$
consequently, for every constant sheaf $F$ on~$X$, the canonical
morphism
$$
F \to i_*i^*F
$$
is an isomorphism. It follows that, if we suppose the
theorem proved for $fi$ and $i$, then, after
restriction of~$S$ to a nonempty open set, $(i^*F,i)$ and $(i^*F,fi)$
are cohomologically proper relative to~$S$ in dimension
$\le 0$. The same therefore holds for~$(F,f)$ (\Ref{XIII.1.6}~2)). Since moreover $f_* F=(fi)_*(i^*F)$ is constructible, this
completes the reduction.

1) 3. End of the proof.

We may assume $S$ normal (EGA~IV 7.8.3). We can find a
compactification of~$X_s$:
% original p. 398
$$
\xymatrix{ X_s \ar[r]^{j_s} \ar[d]_{f_s} & P_s \ar[ld]^{g_s} \\ Y_s }
$$
where $j_s$ is a dominant open immersion and $g_s$ a
proper morphism; after making a radicial extension of~$\kres(s)$ and
replacing $P_s$ by its normalization, which does not change $X_s$, we
may assume $P_s$ geometrically normal. After
restricting $S$ to a nonempty open set and making a surjective
radicial extension if necessary, we may assume that the diagram above
comes from a diagram
$$
\xymatrix{ X \ar[r]^{j} \ar[d]_{f} & P \ar[ld]^{g} \\ Y }
$$
where $P$ is a scheme normal over~$S$, $j$ an open immersion
schematically dominant relative to~$S$ and $g$ a
proper morphism (EGA~IV 6.9.1, 9.9.4 and 9.6.1). For every finite
constant sheaf of sets $F$ on~$X$ with value $C$, $j_*F$ is the
constant sheaf with value $C$ (SGA~4 2.14.1), and the same holds
after any change of base $S' \to S$; it
follows that $(F,j)$ is cohomologically proper relative
to~$S$ in dimension $\le 0$. The same therefore holds for~$(F,f)$,
since $g$ is proper \eqref{XIII.1.8}. Since $g$ is proper
$f_*F=g_*C_P$ is constructible, which completes the
proof of 1)\kern2pt a).

\Subsection*{2) Case of a constructible sheaf of sets}

Let $F$ be a constructible sheaf of sets on~$X$. By
SGA~4 IX 2.14~(ii), we can find a finite family of morphisms
$p_i \colon Z_i \to X$, and, on each $Z_i$, a finite constant sheaf of sets
$C_i$, such that we have a monomorphism
% original p. 399
$$
j \colon F \to \prod_{i} {p_i}_* C_i =G \quoi.
$$
By~1)\kern2pt a),
we may assume, after restricting $S$ to a nonempty
open set if necessary, that the $(C_i,fp_i)$ are cohomologically proper relative
to~$S$ in dimension $\le 0$, and that the $f_*{p_i}_*C_i$ are
constructible. We already conclude from this that $(G,f)$ is
cohomologically proper relative to~$S$ in dimension $\le 0$
\eqref{XIII.1.9}, hence that $(F,f)$ is cohomologically proper
relative to~$S$ in dimension $\le -1$, and that $f_*F$ is
constructible. Let~$K$ be the amalgamated sum $K= G \coprod_{F} G$;
since $F$ and $G$ are constructible, so is~$K$. We
therefore conclude from the preceding that, after restricting
$S$ to a nonempty open set if necessary, we may assume that $(K,f)$ is
cohomologically proper relative to~$S$ in dimension $\le
-1$. It then follows from \Ref{XIII.1.13}~1) that $(F,f)$ is
cohomologically proper relative to~$S$ in dimension $\le 0$.

\Subsection*{3) Case of constant sheaves of groups}

If $F$ is a constant sheaf of groups on~$X$, we denote by $\Phi$ the
stack of torsors under $F$.

3) 1. Let us first show that, after restricting $S$ to a
nonempty open set if necessary, for every constant sheaf of groups $F$ on~$X$,
of order prime to the residue characteristics of~$S$,
$(F,f)$ is cohomologically proper relative to~$S$ in
dimension $\le 0$, and that $f_*\Phi$ is constructible.

For this we reduce to the case where $X$ is smooth over~$S$. After
making a finite extension of~$\kres(s)$ if necessary, which is permissible since one
can regard it as the composite of an \'etale extension and
of a radicial extension, we can find a proper
surjective morphism $p_s \colon X_{1s}\to X_s$, where $X_{1s}$ is a
scheme smooth over~$S$ of the same dimension as $X_s$, and, after
restricting $S$ to a nonempty open set if necessary, we may assume that
$p_s$ comes from a proper surjective morphism $p \colon X_1 \to X$,
where $X_1$ is a scheme smooth over~$S$
% original p. 400
(EGA~IV 9.6.1 and 12.1.6). Let $X_2=X_1 \times_X X_1$, $q \colon
X_2 \to X$ the canonical morphism. We have an exact diagram of stacks
on~$X$
$$
\xymatrix@C=.5cm{\Phi \ar[r] & p_*p^* \Phi \ar@<2pt>[r]\ar@<-2pt>[r] &
q_*q^*\Phi }
$$
(\Ref{XIII.1.11}~2)). The theorem being assumed
proved in the smooth case, we see first of all that we may
assume $f_*p_*p^*\Phi$ constructible; hence so is
$f_*\Phi$ (\Ref{XIII.3.1.1} below). Moreover, by \Ref{XIII.1.6}~2),
we may assume $(p_*p^*\Phi,f)$ cohomologically proper
relative to~$S$ in dimension $\le 0$; it follows that
$(\Phi,f)$ is cohomologically proper relative to~$S$ in
dimension $\le -1$. We may therefore assume that $(q_*q^*\Phi,f)$ is
cohomologically proper relative to~$S$ in dimension $\le -1$,
and this implies that $(\Phi,f)$ is cohomologically proper
relative to~$S$ in dimension $\le 0$ (\Ref{XIII.1.12}~1)).

We then reduce as in 1)\kern2pt2 to the case where $X$ is smooth
and affine over~$S$. Let then
$$
\xymatrix{ X \ar[r]^{i} \ar[d]_{f} & P \ar[ld]^{q} \\ Y }
$$
be a compactification of~$X$, where $i$ is a dominant open
immersion and $q$ a proper morphism. Since we have $\dim P_s=\dim X_s$,
we can apply the hypothesis of resolution of
singularities to $P_{\overline s}$. After making an
\'etale extension and a radicial extension of~$S$ if necessary, we can
find a proper morphism $r \colon Z \to P$, where $Z$ is smooth
over~$S$, $r^{-1}(X)\simeq X$, and where $r^{-1}(X)$ is the
complement in $Z$ of a divisor with normal crossings
relative to~$S$. Every torsor under $F$ is then
tamely ramified on~$Z$ relative to~$S$
(\Ref{XIII.2.3}~b)). It therefore follows from~\Ref{XIII.2.7} that $(F,f)$
is cohomologically proper relative to~$S$ in dimension $\le
0$, which proves our assertion.

3) 2. Reduction to the case where $X$ is smooth over~$S$.
% original p. 401
\par
After making a finite extension of~$\kres(s)$ if necessary, we can find a
proper surjective morphism $p_s\colon X_{1s}\to X$, where~$X_{1s}$
is a scheme smooth over~$s$, and we may assume that $p_s$
comes from a proper surjective morphism $p\colon X_1 \to X$, where
$X_1$ is smooth over~$S$. Suppose the theorem
proved for $fp$ and let us prove it for $f$. Let $F$ be a
finite constant sheaf of groups on~$X$, of order prime to the
residue characteristics of~$S$ and $\Phi$ the stack of
torsors under $F$. Let $X_2=X_1 \times_X X_1$, $X_3=X_1 \times_X
X_1 \times_X X_1$, $q\colon X_2 \to X$, $r\colon X_3 \to X$ the
canonical morphisms. By \Ref{XIII.1.11}~2), we have an
exact diagram of stacks
$$
\xymatrix@C=.5cm{\Phi \ar[r] & p_*p^*\Phi \ar@<2pt>[r]\ar@<-2pt>[r] &
q_*q^*\Phi \ar@<4pt>[r] \ar[r] \ar@<-4pt>[r] & r_*r^*\Phi}\quoi.
$$
To prove that $(\Phi,f)$ is cohomologically proper relative
to~$S$ in dimension $\le 1$, it suffices to show that the same
holds for~$(p_*p^*\Phi,f)$, that $(q_*q^*\Phi, f )$ is
cohomologically proper relative to~$S$ in dimension $\le 0$
and that $(r_*r^*\Phi,f)$ is cohomologically proper relative
to~$S$ in dimension $\le -1$ (\Ref{XIII.1.12}~2)). By~3)\kern.3em 1
above we may assume that, for every finite constant sheaf of
groups~$F$, $(q^*\Phi,fq)$, $(q^*\Phi,q)$, $(r^*\Phi,fr)$, $(r^*\Phi,r)$
are cohomologically proper relative to~$S$ in dimension
$\le 0$. It then follows from \Ref{XIII.1.6}~2) that
$(q_*q^*\Phi,f)$ and $(r_*r^*\Phi,f)$ are cohomologically proper
relative to~$S$ in dimension $\le 0$. The theorem
being assumed proved in the smooth case, $(p^*\Phi,
fp)$ and $(p^*\Phi,p)$, hence also $(p_*p^*\Phi,f)$ (\Ref{XIII.1.6}~2)),
are cohomologically proper relative to~$S$ in dimension
$\le 1$. This indeed shows that $(\Phi,f)$ is cohomologically proper
relative to~$S$ in dimension $\le 1$.

Moreover $f_*p_*p^*\Phi$ is $1$-constructible by hypothesis;
by~\Ref{XIII.3.1} we may assume that $f_*q_*q^*\Phi$ is
constructible; it therefore follows from~\Ref{XIII.3.1.1} below that
$f_*\Phi$ is $1$-constructible.

3) 3. Reduction to the case where $X$ is smooth affine over~$S$.
% original p. 402

By 3)\kern2pt2, we may assume $X$ smooth over~$S$. Let $(U_i)_{i\in
I}$ be a finite covering of~$X$ by affine open sets, and let $X_1$
be the direct sum of the $U_i$, $p\colon X_1 \to X$ the canonical
morphism. Since $p$ is an effective descent morphism for the
category of \'etale sheaves of finite type over
variable schemes, we see as in 3)\kern2pt2 that, if we suppose the
theorem proved for the $U_i$, \ie for $X_1$, it
is also true for~$X$.

3) 4. Case where $X$ is smooth affine over~$S$.

We see as in 3)1 that, after restricting $S$ to a
nonempty open set and making a surjective \'etale extension and a
surjective radicial extension of~$S$ if necessary, we can find a
commutative diagram
$$
\xymatrix{ X \ar[r]^{i} \ar[d]_{f} & P \ar[ld]^{q} \\ Y }
$$
where $P$ is a scheme smooth over~$S$, $X$ the complement
in $P$ of a divisor with normal crossings relative to~$S$ and $g$ a proper
morphism. If $F$ is a constant sheaf
of order prime to the residue characteristics of~$S$, every
torsor under $F$ is tamely ramified on~$P$
relative to~$S$ (\Ref{XIII.2.3}~b)). The fact that $(F,f)$ is
cohomologically proper relative to~$S$ in dimension $\le 1$
and that $f_*\Phi$ is constructible then follows
from~\Ref{XIII.2.7}.
